\documentclass[11pt,a4paper]{article}

\usepackage[T1]{fontenc}
\usepackage[utf8]{inputenc}
\usepackage{lmodern}
\usepackage[a4paper,margin=26mm]{geometry}
\usepackage{amsmath,amssymb,amsthm,mathtools}
\usepackage{booktabs,array}
\usepackage{microtype}
\usepackage{xcolor}
\usepackage{enumitem}
\usepackage{url}
\usepackage[colorlinks=true,linkcolor=blue!55!black,citecolor=blue!55!black,
            urlcolor=blue!55!black]{hyperref}

\newtheorem{theorem}{Theorem}
\newtheorem{lemma}[theorem]{Lemma}
\newtheorem{proposition}[theorem]{Proposition}
\newtheorem{remark}[theorem]{Remark}
\newcommand{\F}{\mathbb F}
\newcommand{\Z}{\mathbb Z}
\newcommand{\bits}{\{0,1\}}
\newcommand{\concat}{\mathbin\Vert}
\newcommand{\sample}{\mathrel{\leftarrow\$}}
\newcommand{\pk}{\ensuremath{\mathsf{pk}}}
\newcommand{\sk}{\ensuremath{\mathsf{sk}}}
\newcommand{\Sign}{\ensuremath{\mathsf{Sign}}}
\newcommand{\Verify}{\ensuremath{\mathsf{Verify}}}
\newcommand{\Hash}{\ensuremath{\mathsf{Hash}}}
\newcommand{\XOF}{\ensuremath{\mathsf{XOF}}}
\newcommand{\DRBG}{\ensuremath{\mathsf{DRBG}}}
\newcommand{\Prb}{\mathop{\rm Pr}}

\title{\bfseries A Two-Signature Complete Pre-Signature DRBG\\
State Rollback Attack on Qing Luan:\\
Equivalent Signing-Witness Recovery and Fresh-Message Forgery\\
Against the Claimed State-Reuse Robustness}
\author{Martin Feussner\\
\small Selmer Center, University of Bergen\\
\small \texttt{martin.feussner@uib.no}}
\date{1 October 2026}

\begin{document}
\maketitle

\begin{center}
\fbox{\parbox{0.92\linewidth}{\small\textbf{Provenance and review status.}
The attack, derivation, implementation, experiments, and initial manuscript
were produced by OpenAI Codex using Daybreak Blue at Ultra reasoning effort
during an AI-run audit of NGCC signature candidates.  A separately instructed
hostile-review agent rebuilt the extraction and forgery, tried to disprove the
claim, and accepted only the conditional complete pre-signature DRBG state
rollback result stated below.  At the date above, no human cryptanalyst has independently verified the
argument, implementation, measurements, or novelty assessment.  This report is
a preliminary disclosure for human review and reproduction.}}
\end{center}
\vspace{0.7em}

\begin{abstract}
Qing Luan is a code-based Fiat--Shamir signature derived from the CROSS R-SDP
protocol.  Its specification separately considers what is effectively
complete pre-signature DRBG state rollback to the same state immediately before
both signing calls on different messages.  It
observes that the per-round seeds repeat while the Fiat--Shamir challenges
change, but nevertheless claims strong robustness under temporary state reuse.

We show that this exact condition yields a practical two-signature witness
recovery.  Suppose one key signs two distinct messages after complete
pre-signature DRBG state rollback to the same state immediately before both
signing calls, so both calls repeat the hidden root seed and the public salt.
Repetition of either value alone is insufficient.  If the
fixed-weight second challenges differ at a round $i$, one
signature's type-1 opening serializes $\mathsf{Seed}_i$, from which
$\eta'_i$ is publicly derived, while the other signature's type-0 opening
serializes
\[
                         v_i=\eta-\eta'_i\pmod 7.
\]
Thus $\eta=\eta'_i+v_i\pmod 7$ is recovered directly.  Because both second
challenges have the same weight, any unequal pair contains such a cross-branch
round.  If the second challenges are identical, a separate type-0/type-0
fallback subtracts two affine responses with unequal first challenges.  The
recovered exponent vector satisfies the public R-SDP relation.  It is an
equivalent signing witness: it need not recover the original secret seed, but
it permits ordinary signatures on arbitrary fresh messages.

A public-input extractor and witness-only signer reproduced the attack at all
four submitted levels.  The affine fallback's main PDF-aligned hostile-review
runs succeeded on $32/32$, $8/8$, $4/4$, and $4/4$ independent keys at 128,
256, 384, and 512 bits.  The combined extractor exercised the simpler
cross-branch route at every level under both PDF-aligned and untouched-source
semantics.  A second DRBG implementation also reproduced the affine route.
Fresh-state and same-message pairs rejected extraction; all corrupted-response,
wrong-key, wrong-message, and forgery-binding controls behaved as required.
For the timed affine runs, mean public extraction time, including verification
of both source signatures, was between 0.010 and 0.313 CPU seconds and peak
memory remained below 5 MiB.

This result is a conditional robustness failure under the complete
pre-signature DRBG state rollback condition stated above.  It does not break
ordinary EUF-CMA signing with fresh
independent randomness.  Its design significance follows from Qing Luan's
explicit robustness statements for essentially this state-reuse condition.  If
those statements are withdrawn, the mechanism is instead a conditional
fault/caller-misuse attack.  Generic Fiat--Shamir nonce-reuse extraction is
well known; the contribution here is limited to Qing Luan's exact
cross-branch seed/correction mechanism, its same-branch affine fallback, the
transport to a complete R-SDP signing witness, and the resulting contradiction
of the submission's explicit robustness claim.
\end{abstract}

\section{Introduction}

Qing Luan supplies four code-based signature parameter sets built from the
restricted syndrome decoding problem over $\F_{127}$ and the CROSS R-SDP
identification protocol~\cite{qingluan-spec,cross-rsdp}.  It executes many
parallel rounds, commits to secret-independent masks, derives two
Fiat--Shamir challenges, and opens each round in one of two forms.  The
construction uses an SM3-based deterministic random-bit generator to sample a
per-signature root and salt.

Randomness reuse is a standard special-soundness hazard for Fiat--Shamir
signatures.  Hedged derivations bind signing randomness to both the secret key
and message precisely because a stateful generator can be reset, cloned, or
restored from a virtual-machine snapshot~\cite{aranha-hedged}.  Qing Luan's
supporting analysis recognizes the relevant event.  Section~3.3.4 describes
what is effectively complete pre-signature DRBG state rollback to the same
state immediately before both signing calls on different messages: the round
seeds repeat while the challenges differ.  It then calls
the construction strongly robust under temporary DRBG state reuse;
Section~3.3.6 again says that the state-reuse mechanism is robust.

The repeated round seed is enough to expose the witness.  The most direct
route uses different second-challenge branches at one round: the type-1
opening supplies the round seed and hence $\eta'_i$, while the type-0 opening
supplies $v_i=\eta-\eta'_i$.  Adding them gives the long-term exponent vector.
If the second challenges happen to be identical, two type-0 affine responses
with different first challenges provide an independent fallback.  The
resulting witness is publicly checkable and can drive the ordinary prover on
any new message.

Our contributions are as follows.
\begin{enumerate}[leftmargin=*,itemsep=2pt]
  \item We give exact public extraction equations for two distinct-message
        signatures produced after complete pre-signature DRBG state rollback
        to the same state immediately before both signing calls, which repeats
        both the hidden root and public salt.
  \item We show that one cross-branch round recovers an equivalent R-SDP
        signing witness.  We retain an affine type-0/type-0 extractor as a
        fallback when the second challenge repeats.
  \item We calculate ideal-model failure probabilities and keep those models,
        including their digest-collision limits, separate from concrete
        serialized-transcript experiments.
  \item We implement a public-input extractor and a signer that receives only
        the recovered witness.  Fresh-message forgeries pass the ordinary
        verifier at all four security levels.
  \item We repair the submitted code to the PDF's independent indexed-leaf and
        one-based domain rules before the main tests, and use untouched source
        only as corroboration.
  \item We delimit the claim as a failure of the explicit robustness
        statements for the complete pre-signature DRBG state rollback condition
        stated above.  We make no ordinary fresh-randomness EUF-CMA claim.
\end{enumerate}

The generic lesson that repeated Fiat--Shamir randomness may expose a witness
is not new.  Qing Luan's own special-soundness discussion describes extraction
from complementary openings.  The candidate-specific point is that the
submission identifies essentially the same complete pre-signature DRBG state
rollback condition and still asserts robustness; its serialized seed and
correction make complete witness recovery immediate.

\section{Qing Luan}

\subsection{Public relation and keys}

All arithmetic in the protocol is over $\F_p$ with
\[
 p=127,\qquad z=7,\qquad g=2,\qquad
 E=\langle g\rangle=\{1,2,4,8,16,32,64\}\subset\F_p^*.
\]
For code length $n$ and dimension $k$, set $r=n-k$ and expand the
public seed into
\[
                  H=[V\mid I_r]\in\F_p^{r\times n}.
\]
Key generation samples $\eta\in\Z_7^n$, sets
\begin{equation}
              e=g^\eta\in E^n,\qquad s=eH^\top,             \label{eq:public}
\end{equation}
and publishes a seed for $H$ together with $s$.  The submitted secret key is a
seed that regenerates $\eta$.

For signing, the seed is only a compact representation.  Any vector
$\widetilde\eta\in\Z_7^n$ satisfying
\begin{equation}
                  g^{\widetilde\eta}H^\top=s               \label{eq:witness}
\end{equation}
is an equivalent signing witness.  It allows the holder to execute every
secret-dependent signing equation, even if it does not invert the original
secret-seed expansion.

The relevant parameters are shown in Table~\ref{tab:params}.  The second
challenge has fixed Hamming weight $w$ among $t$ rounds; therefore exactly
$a=t-w$ rounds use the type-0 response that is central to the attack.

\begin{table}[ht]
\centering
\caption{Qing Luan parameters and serialized sizes.}
\label{tab:params}
\begin{tabular}{lrrrrrrr}
\toprule
profile & $(n,k,r)$ & $t$ & $w$ & $a=t-w$ & $|\pk|$ & $|\sk|$ & $|\sigma|$\\
\midrule
128 & $(127,76,51)$   & 256  & 212 & 44  & 77  & 32  & 18,720\\
256 & $(251,150,101)$ & 512  & 424 & 88  & 153 & 64  & 74,248\\
384 & $(370,221,149)$ & 763  & 631 & 132 & 227 & 96  & 164,940\\
512 & $(491,293,198)$ & 1018 & 842 & 176 & 302 & 128 & 292,816\\
\bottomrule
\end{tabular}
\end{table}

\subsection{Per-round state}

Each signing call draws a hidden root, denoted $\mathsf{Seed}$, and a public
salt.  The written specification derives an independently indexed
$\mathsf{Seed}_i$ for every round.  Expanding it produces
\[
          \eta'_i\in\Z_7^n,\qquad u'_i\in\F_p^n.
\]
The signer forms
\begin{align}
 v_i   &= \eta-\eta'_i \pmod 7,                              \label{eq:v}\\
 e'_i  &= g^{\eta'_i},                                      \label{eq:eprime}\\
 u_i   &= g^{v_i}\odot u'_i,\qquad
 s'_i = u_iH^\top.                                          \label{eq:commitstate}
\end{align}
Here $\odot$ denotes coordinatewise multiplication.  Commitments bind $s'_i$,
$v_i$, the round seed, the salt, and their domains.

After hashing the message and commitments, the first Fiat--Shamir challenge
contains a nonzero scalar
\[
                         \beta_i\in\F_p^*.
\]
The signer computes
\begin{equation}
                         y_i=u'_i+\beta_i e'_i.               \label{eq:y}
\end{equation}
The fixed-weight second challenge selects one of two openings.  A type-1
opening discloses the round seed.  A type-0 opening serializes $y_i$, $v_i$,
and the complementary commitment.  From those public values the verifier
reconstructs
\begin{equation}
        (g^{v_i}\odot y_i)H^\top-\beta_i s=s'_i,              \label{eq:verify}
\end{equation}
because
\[
 g^{v_i}\odot y_i
 =g^{v_i}\odot u'_i+\beta_i g^{v_i+\eta'_i}
 =u_i+\beta_i e.
\]

\subsection{The stated complete pre-signature DRBG state rollback scenario}

The signing API obtains the root and salt from two sequential calls to its
DRBG.  Complete pre-signature DRBG state rollback to the same state immediately
before both signing calls repeats both draws.  For two different messages, the repeated
root and salt reproduce every
$\mathsf{Seed}_i$, $\eta'_i$, $u'_i$, $v_i$, and round commitment.  Message
binding changes the Fiat--Shamir challenges with overwhelming probability.

This is the scenario described in Section~3.3.4 of the submission, using its
``reset'' terminology: round seeds repeat while challenges differ.  The
specification's conclusion of ``strong robustness under DRBG temporary-state
reuse'' overlooks that complete pre-signature DRBG state rollback reproduces
both draws.  The attack below does not arise when the DRBG advances normally
and supplies independent signing state.

\section{The rollback extractor}

\subsection{Cross-branch extraction}

Let $\sigma$ and $\sigma'$ be valid signatures on distinct messages under one
public key.  Assume complete pre-signature DRBG state rollback to the same state
immediately before both signing calls, so both use the same hidden root and
public salt.  Let $B$ and $B'$ be
the respective sets of type-1 positions.  Both sets have cardinality $w$.

Suppose first that $B\ne B'$.  There is an index in their symmetric
difference; without loss of generality, choose $i\in B\setminus B'$.  The
type-1 opening in $\sigma$ serializes $\mathsf{Seed}_i$.  Applying the public
round expansion to that seed, the repeated salt, and the round domain gives
$\eta'_i$.  The type-0 opening at the same index in $\sigma'$ serializes
$v'_i$.  Complete pre-signature DRBG state rollback makes the hidden round
state identical, so
\begin{equation}
                    \widehat\eta=\eta'_i+v'_i=\eta\pmod 7.
                                                        \label{eq:crossbranch}
\end{equation}
One cross-branch position therefore recovers every coordinate of the long-term
exponent witness.  No first-challenge difference or discrete logarithm is
needed.

\subsection{Same-branch affine fallback}

If $B=B'$, every type-0 position is shared.  Consider one such round $i$.
Repetition of the round state gives
\begin{align}
 y_i  &=u'_i+\beta_i e'_i,\nonumber\\
 y'_i &=u'_i+\beta'_i e'_i.                                 \label{eq:twoy}
\end{align}
If $\beta_i\ne\beta'_i$, coordinatewise subtraction and division in
$\F_p$ yield
\begin{equation}
             \widehat e'_i=(y_i-y'_i)(\beta_i-\beta'_i)^{-1}.
                                                                    \label{eq:extracteprime}
\end{equation}
The serialized correction then gives
\begin{equation}
                     \widehat e=g^{v_i}\odot\widehat e'_i.   \label{eq:extracte}
\end{equation}

\begin{theorem}[Conditional witness recovery]
Suppose two honestly generated Qing Luan signatures under the same key sign
distinct messages after complete pre-signature DRBG state rollback to the same
state immediately before both signing calls, repeating both the hidden root and
public salt.  If their
second-challenge vectors differ, Equation~\eqref{eq:crossbranch} recovers the
long-term exponent vector $\eta$.  If those vectors are identical but some
type-0 round has $\beta_i\ne\beta'_i$,
Equations~\eqref{eq:extracteprime} and \eqref{eq:extracte} recover the
long-term restricted error $e$, whose coordinatewise discrete logarithms give
$\eta$.  In either case, $\eta$ is an equivalent signing witness.
\end{theorem}

Repetition of either the hidden root or public salt alone does not satisfy the
theorem's rollback hypothesis and is insufficient for this extraction.

\begin{proof}
For a cross-branch round, the type-1 seed expands to the same $\eta'_i$ used
to compute the other signature's type-0 correction.  Equation~\eqref{eq:v}
then proves Equation~\eqref{eq:crossbranch}.  For the affine fallback,
complete pre-signature DRBG state rollback repeats $\eta'_i$ and $u'_i$.
Subtracting
Equation~\eqref{eq:twoy} cancels $u'_i$ and gives
$(\beta_i-\beta'_i)g^{\eta'_i}$, so
$\widehat e'_i=g^{\eta'_i}$.  By Equation~\eqref{eq:v},
\[
 \widehat e=g^{\eta-\eta'_i}\odot g^{\eta'_i}=g^\eta=e.
\]
Every coordinate lies in the seven-element group $E$, so its discrete
logarithm to base $g$ is found by a seven-entry table.  The resulting exponent
vector satisfies Equation~\eqref{eq:witness}.
\end{proof}

The public extractor first verifies both input signatures and derives both
challenge vectors.  It tries every cross-branch position using the serialized
type-1 seed and opposite type-0 correction; if necessary, it tries shared
type-0 positions with unequal first challenges.  It accepts a candidate only
if every exponent is in $\Z_7$, the corresponding restricted vector satisfies
$eH^\top=s$, and every other usable position recovers that same vector.  Its
interface is
\[
       (\pk,M,\sigma,M',\sigma')\longmapsto\eta\ \text{or failure}.
\]
It receives no secret key, secret-seed expansion, DRBG state, expected witness,
hidden root, unexposed round seed, or signer instrumentation.  Its only round
seed input is the one already serialized by a public type-1 opening.

\subsection{Fresh-message forgery}

Given the recovered $\eta$, the adversary can choose fresh signing randomness
and compute the signing state using Equations~\eqref{eq:v}--\eqref{eq:y}.
It can then serialize an ordinary proof for any message.  The secret seed is
unnecessary.  The accompanying
witness-only signer takes $(\pk,M,\eta)$ and produces a conventional Qing Luan
signature.  Thus the result is equivalent signing-witness recovery followed by
fresh-message forgery, not merely a distinguisher or transcript correlation.

\subsection{Ideal-model success probability}

Let $B,B'\subseteq\{1,\ldots,t\}$ be the type-1 positions, with
$|B|=|B'|=w$.  In the ideal model where both fixed-weight challenges are
independently uniform, cross-branch extraction fails exactly when $B=B'$:
\begin{equation}
       P_{\rm cross\text{-}fail}^{\rm ideal}=\binom{t}{w}^{-1}.
                                                        \label{eq:crossfailure}
\end{equation}
If that collision occurs, all $a=t-w$ type-0 positions are shared.  An
additional abstraction that treats their first-challenge scalars as independent
uniform elements of $\F_p^*$ makes the affine fallback fail with conditional
probability $126^{-a}$.  Ignoring the finite digest layers would therefore give
the product $(\binom{t}{w}126^a)^{-1}$.

\begin{table}[ht]
\centering
\caption{Cross-branch failure in the uniform fixed-weight ideal model.  These
are not exact probabilities for conditioned serialized signatures.}
\label{tab:prob}
\begin{tabular}{lrr}
\toprule
profile & $a=t-w$ & $-\log_2 P_{\rm cross\text{-}fail}^{\rm ideal}$
        \\
\midrule
128 & 44  & 165.543619\\
256 & 88  & 334.510906\\
384 & 132 & 502.320455\\
512 & 176 & 671.306477\\
\bottomrule
\end{tabular}
\end{table}

Equation~\eqref{eq:crossfailure} is exact only for its stated uniform
fixed-weight model.  A layered ideal-hash calculation also shows why the
abstract product must not be reported as the concrete failure probability.
Let $h=2\lambda$, $q=2^{-h}$, $D=\binom{t}{w}$, and
$d=q+(1-q)q$.  Here $d$ models a collision in the message digest or,
otherwise, in $\mathsf{digest\_chall1}$; either event repeats the entire
challenge cascade.  Conditional on avoiding it, equality of the second
challenge has modeled probability $c=q+(1-q)/D$, accounting first for a
collision in its $h$-bit digest and then for the fixed-weight sampler.  The
layered model gives
\[
 P_{\rm cross\text{-}fail}=d+(1-d)c,\qquad
 P_{\rm combined\text{-}fail}=d+(1-d)126^{-a}c.
\]
The combined failure is consequently dominated by about $2^{-(2\lambda-1)}$,
not by the much smaller digest-free product.  These remain idealized
calculations; concrete challenges are hashes of conditioned serialized
transcripts.  We report real-signer results separately.

\section{Specification alignment and implementation}

\subsection{Main PDF-aligned experiment}

The main experiment starts from fresh copies of the submitted source and makes
two semantic repairs required by the PDF:
\begin{enumerate}[leftmargin=*,itemsep=2pt]
  \item derive each round leaf independently as an indexed expansion, rather
        than slicing one unindexed XOF stream;
  \item use the PDF's one-based round number in every $i+c$ domain, rather than
        the source's zero-based value.
\end{enumerate}
The bundled protocol reference's key-expansion constants are retained where
the main PDF abbreviates them.  The complete patch accompanies the reproducer.
The attack is then tested against the repaired signer and verifier, so neither
submitted deviation is an attack premise.

For each trial, the producer generates a fresh key and distinct messages,
induces complete pre-signature DRBG state rollback to the same state immediately
before both signing calls, and outputs two valid serialized signatures.  The extractor receives
only the public tuple.  A separate replay check confirms that the two
sequential DRBG outputs---the hidden root and serialized salt---repeat exactly,
but does not pass them to the extractor.

The witness-only signer receives the recovered $\eta$ and an unqueried message.
The ordinary verifier checks the resulting signature.  AddressSanitizer and
UndefinedBehaviorSanitizer runs cover the PDF-aligned 128- and 512-bit paths.

\subsection{Untouched-source corroboration}

The same independent harness is also linked to byte-for-byte untouched
submitted core copies.  This corroboration shows that the effect was not
introduced by the conformance patch.  It is not used to replace the normative
PDF-aligned result.  A second experiment uses a standalone implementation of
the documented DRBG, which rules out an unrelated rotate operation in the
submitted API generator as the cause.

\section{Results}

Table~\ref{tab:results} gives the independently rebuilt affine hostile-review
runs and the combined cross-branch reproductions.  Every numerator is a
complete sequence of two valid signatures produced after complete
pre-signature DRBG state rollback to the same state immediately before both
signing calls, public witness recovery, an unqueried-message
signature using only that witness, ordinary-verifier acceptance, and all
negative controls.

\begin{table}[ht]
\centering
\small
\caption{End-to-end recovery and fresh-message forgery.}
\label{tab:results}
\begin{tabular}{lrrrr}
\toprule
semantics and randomness backend & 128 & 256 & 384 & 512\\
\midrule
affine: PDF-aligned / submitted API DRBG & 32/32 & 8/8 & 4/4 & 4/4\\
affine: PDF-aligned / standalone documented DRBG & 4/4 & 4/4 & 2/2 & 2/2\\
affine: untouched submitted core / API DRBG & 4/4 & 4/4 & 2/2 & 2/2\\
combined: PDF-aligned / documented DRBG & 1/1 & 1/1 & 1/1 & 1/1\\
combined: untouched core / documented DRBG & 1/1 & 1/1 & 1/1 & 1/1\\
\bottomrule
\end{tabular}
\end{table}

In the main PDF-aligned/API affine runs, the observed minimum, mean, and maximum
type-0 overlaps were
\[
 3/6.844/12,\quad 12/16.375/21,\quad
 18/22.250/28,\quad 22/33.000/49
\]
from 128 through 512.  Occasional overlaps with equal $\beta$ were skipped; all
remaining candidate rounds recovered the same witness.

For those affine runs, mean extraction CPU time, including two ordinary
signature verifications, was
0.010, 0.052, 0.148, and 0.313 seconds.  Peak resident memory for complete
batches stayed below 5 MiB.  Public algebra costs
$O(tn+rk)$ small-field operations and comparable storage.  Online data is two
signatures: 37,440, 148,496, 329,880, and 585,632 bytes at the four levels,
plus the public key and messages.

Every trial required the following controls to pass:
\begin{itemize}[leftmargin=*,itemsep=2pt]
  \item signatures generated from different advancing DRBG states verified,
        but extraction rejected;
  \item same-message complete pre-signature DRBG state rollback to the same
        state immediately before both signing calls produced identical
        signatures and no unequal first challenge, and the public
        attack interface rejected equal messages;
  \item a one-bit corruption in a serialized type-0 $y_i$ made verification
        and extraction reject;
  \item a separately generated public key and an altered source message made
        verification and extraction reject;
  \item the valid fresh-message forgery failed when paired with a previously
        queried message.
\end{itemize}

\section{Scope, prior work, and novelty}

The attack requires all of the following:
\begin{enumerate}[leftmargin=*,itemsep=2pt]
  \item one long-term Qing Luan key;
  \item two distinct messages;
  \item complete pre-signature DRBG state rollback to the same state
        immediately before both signing calls, repeating both the hidden root
        and public salt;
  \item either unequal second challenges, which guarantee a cross-branch
        round, or---for the affine fallback after an identical second
        challenge---one shared type-0 round with unequal first challenges.
\end{enumerate}
It does not apply to an ordinary signing oracle whose DRBG state advances and
supplies fresh independent values.  Repetition of either the hidden root or
public salt alone is insufficient.

Randomness-reuse extraction in Fiat--Shamir signatures is established prior
art.  Aranha, Orlandi, Takahashi, and Zaverucha analyze fault resilience and
hedged derivations that bind the secret key, message, and external
randomness~\cite{aranha-hedged}.  ZKFault studies fault attacks on compressed
zero-knowledge signatures including CROSS through a different fault
surface~\cite{zkfault}.  The Qing Luan specification itself describes
special-soundness extraction from complementary openings and predicts repeated
round seeds under what is effectively the complete pre-signature DRBG state
rollback scenario.

Accordingly, we claim neither the generic nonce-reuse idea nor a new
special-soundness principle.  The narrow contribution is the
Qing-Luan-specific chain
\[
 (\mathsf{Seed}_i,v'_i)\longmapsto(\eta'_i,v'_i)
        \longmapsto\eta,
\]
together with the same-branch affine fallback, its public-input end-to-end
validation at every parameter level, and the fact that it directly contradicts
the submission's explicit robustness conclusion for that rollback scenario.
A final best-effort check at 16:28~UTC on 1 October 2026 found that the Qing
Luan page on ngcc.dev still listed only \texttt{sign-20-1}, an unrelated Medium
proof gap in the 128-bit parameter's quantum accounting~\cite{ngcc-qingluan}.
No public report applying this exact Qing Luan transcript pair to complete
pre-signature DRBG state rollback was found; that negative search is not a
definitive novelty proof.

If the specification withdraws its robustness statements and treats correct
fresh DRBG evolution as an unconditional environmental requirement, the result
should be classified as a conditional fault/caller-misuse attack.  It should
still inform implementations that support VM snapshots, deterministic test
replay, external DRBG injection, or recoverable persistent state.

\section{Repair}

A robust repair should hedge the per-signature root with both secret and
message-bound material.  For example, derive separate root and salt subkeys
from a construction of the form
\[
 \mathsf{PRF}(\mathsf{domain}\concat\sk\concat\Hash(\pk)
              \concat\Hash(M)\concat\mathsf{external\_randomness}).
\]
Binding the public-key and message digests into every indexed round-seed
derivation prevents a repeated external-randomness value from reproducing the
same masks across different messages.  Root and salt derivations should use
separate domain labels.  A persistent anti-rollback counter may provide
defense in depth, but the cryptographic derivation should remain safe if that
state repeats.

The exact repair requires a revised byte-level specification and proof.
Alternatively, the current text must remove the reset/state-reuse robustness
claims and state that complete pre-signature DRBG state rollback to the same
state immediately before both signing calls on different messages exposes the
signing witness.

\section{Conclusion}

Two Qing Luan signatures created after complete pre-signature DRBG state
rollback to the same state immediately before both signing calls, which repeats
both the hidden root and public salt, expose the
long-term exponent witness through one complementary opening; a repeated
type-0 round supplies an independent affine fallback.  The recovered exponent
vector is an equivalent signing witness and produces ordinary fresh-message
signatures.  Repetition of either the hidden root or public salt alone is
insufficient.  The attack is polynomial, has near-unit success in the stated
ideal challenge model, and reproduces at every submitted security level.

The result has a strict boundary: it is not an ordinary fresh-randomness
EUF-CMA break.  Its design significance follows from Qing Luan's explicit
robustness claim for essentially this condition: complete pre-signature DRBG
state rollback to the same state immediately before both signing calls.  That
condition creates the complementary public openings.  If the claim is
withdrawn, the result should instead be classified as a conditional
fault/caller-misuse attack.

\section*{Acknowledgements}

The computations were performed on the Norwegian Research and Education Cloud
(NREC), using resources provided by the University of Bergen and the University
of Oslo.

\section*{Artifact availability}

The accompanying reproducer contains the PDF-alignment patch, independent
public-input extractor, witness-only signer, untouched-source corroboration,
negative controls, exact ideal-probability calculation, and one-command build
and run instructions.  It does not require either source secret key as an
extractor input.

\begin{thebibliography}{9}

\bibitem{qingluan-spec}
Wu Guangfu, Liu Lei, Wu Rui, Liu Lulu, and Li Kangjun,
\newblock \emph{Qingluan Post-Quantum Digital Signature Algorithm Text},
\newblock NGCC submission specification, 2026.

\bibitem{cross-rsdp}
M.~Baldi, M.~Battaglioni, F.~Chiaraluce, A.-L.~Horlemann,
E.~Persichetti, P.~Santini, and V.~Weger,
\newblock ``A new path to code-based signatures via identification schemes
with restricted errors,''
\newblock \emph{Advances in Mathematics of Communications}, vol.~19, no.~5,
pp.~1360--1381, 2025.

\bibitem{aranha-hedged}
D.~F. Aranha, C.~Orlandi, A.~Takahashi, and G.~Zaverucha,
\newblock ``Security of hedged Fiat--Shamir signatures under fault attacks,''
\newblock in \emph{Advances in Cryptology---EUROCRYPT 2020}; full version,
IACR Cryptology ePrint Archive, Report 2019/956, 2021.

\bibitem{zkfault}
P.~Mondal, S.~Adhikary, S.~Kundu, and A.~Karmakar,
\newblock ``ZKFault: Fault attack analysis on zero-knowledge based
post-quantum digital signature schemes,''
\newblock in \emph{Advances in Cryptology---ASIACRYPT 2024}, pp.~132--167;
IACR Cryptology ePrint Archive, Report 2024/1422, 2024.

\bibitem{ngcc-qingluan}
NGCC Security Reports,
\newblock \emph{Qing Luan (sign-20)},
\newblock \url{https://ngcc.dev/reports/sign-20.html}, accessed 1 October 2026.

\end{thebibliography}

\end{document}
